\chapter{Reichenbach: probability logic, posits, and inductive commitments}
\label{sec:p15-reichenbach-probability-logic}

Reichenbach's probability logic is not ordinary Bayesian updating with
an unspecified prior. Its central objects are coordinated sequences of
events and propositions; its numerical values have a frequency
interpretation; its inductive application requires a posit about an
unobserved continuation. Formal calculation, empirical coordination,
and the proposed vindication of induction are separate steps.

\section{Source presentations and the task of coordination}
\label{sec:p15-reichenbach-presentations}

The directly read original presentations here are the French paper
\emph{Les fondements logiques du calcul des probabilités}, catalogued
in 1937, and the weight/probability-logic discussion in an English
\emph{Experience and Prediction} chapter excerpt, a work originally
published in 1938. They are attributed separately. The German
\emph{Wahrscheinlichkeitslehre} (1935) and the revised English
\emph{Theory of Probability} (1949) are distinct versions, not silently
supplied by either original or by a historical quotation. The separate
Eberhardt--Glymour chapter is a later historical reconstruction
\cite{Reichenbach1937Fondements,Reichenbach1938EPChapterReading,
EberhardtGlymour2011ReichenbachReading}.

The French original distinguishes \emph{internal} relations of an
axiomatic mathematical calculus from its \emph{external} application
to physical events. Deduction within the calculus is exact even when
its application concerns uncertainty. A coordination specifies which
events the symbols describe, how two event sequences are paired, and
in which order they are enumerated. It must satisfy the proposed
mathematical relations; choosing a suggestive interpretation is not
itself proof of that correspondence
\cite[chap.~I; chap.~II, secs.~1--2 and6]{Reichenbach1937Fondements}.

The historical authors relate this stance to Reichenbach's changing
views about relativity, causation, and quantum uncertainty. Their
chronology distinguishes the early neo-Kantian thesis of lawful
distribution from the later abandonment of an a priori guarantee of
convergence. Their account of his interpretation of Heisenberg is an
intellectual-historical claim, not a theorem establishing the quantum
theory or a universal metaphysical interpretation of it.

\section{The original relational language and its calculus}
\label{sec:p15-reichenbach-relational-calculus}

\begin{hbdefinition}{def:p15-reichenbach-probability-implication}{Probability implication and reference sequence}
In the French original, paired sequences \(x_i,y_i\) are coordinated
in a stipulated order. Classes \(A,B\) describe event types. The
probability implication, rendered here as
\(A\mathrel{\underset{p}{\supset}}B\), abbreviates the assertion
for the paired sequences that membership of \(x_i\) in \(A\)
implies membership of \(y_i\) in \(B\) with degree \(p\).
The source's functional notation is \(P(A,B)=p\), with reference
class \(A\) \emph{first}. It is not \(P(A\land B)\).

Under the frequency interpretation define, when its denominator is
nonzero,
\[
 F_n(A,B)\defeq
 \frac{\#\{i\leq n:x_i\in A,\ y_i\in B\}}
      {\#\{i\leq n:x_i\in A\}},\qquad
 P(A,B)\defeq\lim_{n\to\infty}F_n(A,B)
\]
where that limit exists. Selecting the indices satisfying the reference
condition gives the corresponding one-term notation on a restricted
sequence. The two-term notation makes the reference restriction
explicit. The pairing need not be identity; a specified temporal
or other coordination can connect \(x_i\) with \(y_i\).
\end{hbdefinition}

The source's implicative schemata have four headings: univocality,
normalization, addition, and multiplication. Univocality says that two
different degrees for the same probability implication can coexist only
for an empty antecedent class; its displayed biconditional explicitly
handles that empty-class case. Normalization connects formal implication
to degree one and requires nonnegative degrees. Addition concerns
consequents exclusive relative to the antecedent. Multiplication connects
\(A\) to \(B\), then \(A.B\) to \(C\), and determines
\(A\) to \(B.C\) by the product of their degrees. The source's
functional form on p.~279 is
\[
\begin{aligned}
 P(A,A\lor B)&=1,& P(A,B.\bar B)&=0,&P(A,B)&\geq0,\\
 P(A,B\lor C)&=P(A,B)+P(A,C)-P(A,B.C),\\
 P(A,B.C)&=P(A,B)P(A.B,C).
\end{aligned}
\]
Here a bar within a class expression denotes complement. The original
empty-class predicate in parentheses has a different typed role from
that complement operation. The later historical chapter's reported
axioms add explicit nonempty-antecedent and existential-degree clauses;
that surface is not presented as a literal transcription of the 1937
schemata.

The original also states a \emph{rule of existence}: an unknown
probability determined as a continuous function of given probability
values exists. Its qualification is continuity at those given values,
not permission to divide by zero or to postulate every desired joint
distribution. The rule permits solving, for instance, an addition
equation for the missing degree. This is a mathematical existence
rule inside the proposed calculus, distinct from the later inductive
rule about empirical continuations
\cite[pp.~279--283 and286--288]{Reichenbach1937Fondements}.

\begin{hbtheorem}{thm:p15-reichenbach-counting-calculus}{proved=annex:p15.reichenbach-counting-calculus}{Counting realization and qualified elimination}
For finite coordinated data, use the uniform measure on its indices
and regard \(A,B,C\) as index events. On every nonempty reference
event, the functional equations above hold. In particular
\[
 P(A,B)+P(A,\bar B)=1.
\]
If both \(A.B\) and \(A.\bar B\) are nonempty,
\[
 P(A,C)=P(A,B)P(A.B,C)+P(A,\bar B)P(A.\bar B,C).
\]
If one branch is empty, its contribution is zero; its conditional
degree is not evaluated. For positive \(P(A,B)\) and \(P(A,C)\),
\[
 P(A.C,B)=\frac{P(A,B)P(A.B,C)}{P(A,C)}.
\]
If all relevant finite frequencies have limits and the denominator
has a positive limit, these identities pass to the limiting
frequencies. These are counting soundness statements on their stated
domains, not a completeness or representation theorem for the entire
historical syntax.
\end{hbtheorem}

The elimination formula reconstructs the method in original §3:
resolve the intermediate class into \(B\) and its complement,
use disjoint addition, then multiply along each branch. Independence
is the additional condition \(P(A.B,C)=P(A,C)\), not an automatic
consequence of the syntax. Bayes' converse calculation is likewise
qualified; probability implication does not license contraposition
in the manner of deductive material implication.

\section{Probability logic of propositions and of propositional series}
\label{sec:p15-reichenbach-series-logic}

In the English source, §35 first assigns symbols \(a,b,\ldots\)
values \(P(a)\in[0,1]\), then equips them with the customary
Boolean connections. One interpretation is the \emph{logic of weights}
of propositions; another is the \emph{logic of propositional series}.
For the latter, \((a_i)\) denotes a sequence obtained from one
propositional function on its event domain, and a connection is made
index by index:
\[
 [(a_i)\lor(b_i)]\defeq(a_i\lor b_i),\qquad
 [(a_i).(b_i)]\defeq(a_i.b_i),
\]
with analogous pointwise negation, implication, and equivalence.
The French original §17 additionally identifies the comma operation
as selection: retain \(b_i\) precisely where \(a_i\) is true,
and renumber the resulting subsequence. This is the frequency
interpretation of the relative degree \(P(a,b)\).

\begin{hbtheorem}{thm:p15-reichenbach-coupling}{proved=annex:p15.reichenbach-coupling}{The actual three-parameter table and its bounds}
For a probability model, put \(u=P(a)\), \(v=P(b)\), and
\(w=P(a.b)\). The possible parameters are exactly
\[
 0\leq u,v\leq1,\qquad
 \max(0,u+v-1)\leq w\leq\min(u,v).
\]
They determine
\[
\begin{array}{c|c}
 \text{connection}&\text{value}\\\hline
 \bar a&1-u\\
 a.b&w\\
 a\lor b&u+v-w\\
 a\supset b&1-u+w\\
 a\equiv b&1-u-v+2w.
\end{array}
\]
If \(u>0\), one may instead specify the coupling
\(r=P(a,b)=w/u\), subject to the corresponding bounds on
\(ur\). The value of the material implication \(a\supset b\)
is \(1-u+ur\), \emph{not} \(r\). Restricting \(u,v\)
to zero and one forces \(w=uv\), giving the classical two-valued
tables for these connections.
\end{hbtheorem}

The English original's worked comparison has \(P(a)=P(b)=1/6\).
If the events are the two different faces of one die, their product
is impossible and their disjunction has probability \(1/3\).
If they are one specified face on each of two independent dice,
the product has probability \(1/36\) and the disjunction has
probability \(11/36\). The same marginal values give different
compound values. This is Reichenbach's point, not a declaration that
all empirical events are independent
\cite[sec.~35, pp.~320--325]{Reichenbach1938EPChapterReading}.

For another calculation, \(u=2/5,v=3/5,w=3/10\) gives coupling
\(3/4\), disjunction \(7/10\), material implication \(9/10\),
and equivalence \(3/5\). The coupling value and the value of the
material conditional must not be conflated.

The formal construction is consequently not truth-functional in the
ordinary many-valued sense: the two argument values alone do not fix
a compound value. The historical authors explain the criticism by
Nagel, Tarski, and Russell using this missing coupling information.
They also identify a different issue: the arithmetic/existential
language of probability degrees and the propositional/sequential
language are not supplied with a fully explicit common recursive
syntax. In particular a nested probability implication needs a
specified type and a semantics for its antecedent; it is not supplied
merely by drawing a second probability arrow.

\section{Sequence order, normality, and missing limits}
\label{sec:p15-reichenbach-normality}

In original §6 the elementary frequency calculus requires the relevant
limits, not a universal irregularity postulate. Original §§7--11 then
classify sequences by their order structure. A \emph{phase probability}
refers to specified predecessor/successor positions. Lack of after-effect
requires that conditioning on a finite predecessor pattern not change
the appropriate later degree. Repeated multiplication then gives
products of phase degrees for those patterns. This is an additional
sequence property, not a theorem that all observations are independent.

The original also distinguishes overlapping counting from counting in
nonoverlapping sections. A property verified with one counting method
need not hold for the other. Its regular divisions select indices in
an arithmetic progression of fixed length \(\lambda\) and phase
\(k\). The \emph{domain of invariance} preserves both the principal
degree and the specified phase probabilities. A \emph{normal} sequence
is without after-effect and has these regular divisions in that domain.
The original uses these conditions to obtain the multiplication and
Bernoulli results in both counting regimes; it credits Copeland's earlier
construction of such sequences. That cited construction is not supplied
by calling a sequence normal
\cite[pp.~288--295]{Reichenbach1937Fondements}.

This original distinction between an elementary calculus and special
ordered subclasses must remain visible beside the historical chapter's
summary in terms of infinite normal sequences. Neither is a license to
replace normality by an undefined assumption of ``randomness''.

\begin{hbtheorem}{thm:p15-reichenbach-frequency-boundaries}{proved=annex:p15.reichenbach-frequency-boundaries}{Three concrete boundaries of limiting-frequency semantics}
The following constructions are possible:
\begin{enumerate}
\item Two binary sequences each have frequency limit \(1/2\), but
      their pointwise conjunction and disjunction have no frequency
      limits. Thus marginal limit existence alone does not make a
      domain closed under Boolean combinations.
\item A reference subsequence has global density zero and is infinite,
      while the frequency of a selected outcome \emph{within} it tends
      to \(1/2\). A quotient of its two zero global densities cannot
      recover that selected-sequence degree.
\item A propositional series has truth frequency one but contains a
      false proposition. Frequency one is therefore not the assertion
      that every proposition in the series is true.
\end{enumerate}
The first construction concerns bare marginal limit existence, not an
already established family with the additional joint/normality conditions.
\end{hbtheorem}

The second construction explains why the counting-calculus theorem
requires a positive limiting denominator when it passes an ordinary
conditional ratio to a limit. Direct subsequence selection is a different
operation. The third is also why the classical zero/one table subcase
must not be taken as a proof of a universal scientific law. The later
authors discuss both closure difficulties and the distinction between
a high frequency of successful consequences and the truth of the
theory itself.

\section{The original modal and two-valued transformations}
\label{sec:p15-reichenbach-modal-transformations}

The French source's own table uses \(p=P(\varphi x_i)\),
\(q=P(\psi y_i)\), and the coupling
\(u=P(\varphi x_i,\psi y_i)\). Its entries are
\(1-p\) for negation, \(p+q-pu\) for disjunction, \(pu\)
for product, \(1-p+pu\) for material implication,
\(1-p-q+2pu\) for equivalence, and \(pu/q\) for reverse
selection when \(q>0\). At \(p=1\), coupling is \(q\);
at \(p=0\), it is indeterminate in the table. The preceding
theorem uses a joint-mass parameter only as an explicitly declared
mathematical reconstruction of these entries, not a change to the
source's notation or an assertion of a total null-reference conditional.

\begin{hbtheorem}{thm:p15-reichenbach-modal-division}{proved=annex:p15.reichenbach-modal-division}{Extensional modalities and approximate threshold division}
On a nonempty index domain, let \(N\) mean that every member of
a propositional series is true, \(\mathcal J\) that every member
is false, and \(\mathrm{Pos}\) that both kinds occur. This is
the source's extensional possibility convention, excluding necessity.
Negation exchanges \(N,\mathcal J\) and preserves Pos. The
complete binary modal table across nonempty domains is
\[
\begin{array}{cc|cccc}
 a&b&a.b&a\lor b&a\supset b&a\equiv b\\\hline
 N&N&N&N&N&N\\
 N&\mathrm{Pos}&\mathrm{Pos}&N&\mathrm{Pos}&\mathrm{Pos}\\
 N&\mathcal J&\mathcal J&N&\mathcal J&\mathcal J\\
 \mathrm{Pos}&N&\mathrm{Pos}&N&N&\mathrm{Pos}\\
 \mathrm{Pos}&\mathrm{Pos}&\mathcal J/\mathrm{Pos}&\mathrm{Pos}/N&
           \mathrm{Pos}/N&\mathcal J/\mathrm{Pos}/N\\
 \mathrm{Pos}&\mathcal J&\mathcal J&\mathrm{Pos}&\mathrm{Pos}&\mathrm{Pos}\\
 \mathcal J&N&\mathcal J&N&N&\mathcal J\\
 \mathcal J&\mathrm{Pos}&\mathcal J&\mathrm{Pos}&N&\mathrm{Pos}\\
 \mathcal J&\mathcal J&\mathcal J&\mathcal J&N&N.
\end{array}
\]
On a finite domain these three categories correspond to frequency
one, strictly between zero and one, and zero. On an infinite domain
those frequency characterizations need not hold.

Classifying every degree above a fixed \(0<\tau<1\) as ``true''
does not preserve conjunction: two independent propositions can each
have degree above \(\tau\) while their conjunction has degree below
\(\tau\). For example degrees \(0.91,0.91\) exceed \(0.9\),
but their product is \(0.8281\).
\end{hbtheorem}

Original §18 calls this threshold/trichotomy procedure a passage by
\emph{division} and explicitly treats it as approximate. The other
passage represents weights through frequencies of underlying two-valued
propositions. It is exact as a mathematical representation only when
the counted propositions and the frequency statement are genuinely
two-valued; Reichenbach says their physical application instead
involves further posits and approximation. His claim of a probability
logic of science is not established by a threshold table alone.

\section{Weights and posits are commitments, not certified truths}
\label{sec:p15-reichenbach-posits}

The English original §34 describes a posit through a wager: one acts
on a statement without thereby claiming to know its truth. Its
\emph{weight} is the predictive contribution which a degree of
probability makes when applied to a single case. The source notes
that its German term \emph{Setzung} has both posit and wager
significations. The French original uses \emph{mise} and
\emph{poids}; those terms are recorded as its own presentation,
not as a new single-case probability hidden in the frequency definition.

The French §19 expressly distinguishes two justifications. An
\emph{appraised} wager has an estimated weight from an appropriate
statistical class. A \emph{blind} wager about a frequency limit is
defended as an approximation method even though its current accuracy
and distance from eventual stability are unknown. Blind does not mean
that an arbitrary numerical probability of its correctness has already
been proved. At the highest level of his proposed network, blind
posits remain; their uncertainty is not eliminated by renaming them.

Reference-class selection is consequently part of the inferential task.
The English source proposes the narrowest class for which reliable
statistics are available, and acknowledges that the resulting weight
depends on knowledge as well as on the event. Its concrete improvement
argument divides a class with success rate \(1/2\) into two equally
represented classes with rates \(1/4\) and \(3/4\). Always predicting
success in the undivided class wins half the cases. Predicting failure
in the first subclass and success in the second wins three quarters.
This argument presupposes those class frequencies and the correct
assignment of the new case; it does not determine the appropriate class
from its name alone.

\begin{hbtheorem}{thm:p15-reichenbach-order-and-appraisal}{proved=annex:p15.reichenbach-order-and-appraisal}{Appraisal improvement and the role of enumeration}
For the stated equally represented two-subclass model, the refined
prediction has success rate \(3/4\), compared with \(1/2\) for
an unrefined constant prediction. More generally, on a finite or
frequency-weighted partition with known rates \(p_j\), selecting
the more frequent binary outcome in each class has success rate
\(\sum_j\alpha_j\max(p_j,1-p_j)\), at least the rate of a
common constant prediction.

If a countably infinite class contains infinitely many members of both
kinds, different enumerations can give any prescribed frequency limit
\(p\in[0,1]\), or no frequency limit. A limiting-frequency
probability therefore requires an order specification, not just an
unordered class.
\end{hbtheorem}

Russell's order objection is thus not answered by saying that the
underlying class is objective. Reichenbach acknowledges order as part
of the specification and regards that dependence as important rather
than fatal. The disagreement is about the intended logical and empirical
role of the order, not the elementary fact that order can matter.
For single historical claims such as Caesar's visit to Britain, the
original appeals to the reliability of chronicles of an appropriate
type. The historical authors question both how that class is determined
and how silence is distinguished from an explicit contrary report.

\section{The straight rule and what its convergence proves}
\label{sec:p15-reichenbach-straight-rule}

The original French §19 states the rule in the form of a wager:
after \(n\) terms with observed frequency \(f^n\), posit that
the continued frequency approaches a limit within \(f^n\pm\delta\).
The historical chapter quotes the 1949 formulation with its additional
qualification: nothing is known about the second-level probability for
a particular limit. That 1949 wording is reported through the
historical source, not claimed to have been read in an acquired 1949
book. Repeated application uses the current empirical frequency;
it is the \emph{straight rule}, not a selected modern prior.

\begin{hbtheorem}{thm:p15-reichenbach-straight-rule-boundary}{proved=annex:p15.reichenbach-straight-rule-boundary}{Finite attainability without a uniform finite certificate}
Let a binary sequence have frequencies
\(f_n=n^{-1}\sum_{i=1}^n x_i\) and suppose \(f_n\to p\).
Then for every \(\varepsilon>0\) there is a finite \(N\)
such that \(|f_n-p|<\varepsilon\) for every \(n\geq N\).
This is conditional pointwise attainability. It supplies neither a
known \(N\) from a finite prefix nor a common such \(N\) for
all convergent sequences.

Every finite binary prefix has continuations with each possible limit
\(p\in[0,1]\). Consequently no finite observed prefix alone
certifies the limit or a nontrivial uniform error bound over that class.
Moreover \(\operatorname{clip}_{[0,1]}(f_n+d_n)\) has the same
limit whenever \(d_n\to0\), and arbitrary predictions through
any fixed initial stage followed by \(f_n\) have the same eventual
limit. Convergence alone does not uniquely select the straight rule.
\end{hbtheorem}

Reichenbach accepts Hume's objections to a deductive or circularly
empirical demonstration that induction succeeds. His proposed change
of task is pragmatic: justify a means to the epistemic end when that
end is attainable, not prove the conclusion true or probable from no
inductive premises. The original explicitly discusses a clairvoyant
who could predict a sequence with no frequency limit. Its response is
that the person's reliability would itself be assessed through a
success/failure classification and induction. It therefore does not
postulate that every raw natural sequence must converge; it appeals
to the availability of suitable classifications, or practical
near-convergence, for predictive success
\cite[pp.~334--341]{Reichenbach1937Fondements}.

The argument about the necessity of such classifications is a
philosophical vindication with an interpretation of \emph{prediction}.
It is not the proven convergence theorem above. The original also
acknowledges the practical-lifetime issue: a mathematical limit reached
only after inaccessible lengths is no useful guarantee. Its claim that
induction is the only known primary route to checking alternatives is
not a theorem quantifying over every possible prediction procedure.

\section{Cross-inductions and the grid of probabilities}
\label{sec:p15-reichenbach-cross-inductions}

Original §19 arranges several sequences of observations in rows of
a grid. Each row gives a blind estimate of its own limit. Across rows,
one counts estimates lying in a stipulated interval and makes a
second-level posit about that proportion. Under further Bernoulli-type
conditions, a Bayes calculation is intended to appraise the original
row posit. The additional conditions and the correctness of the
classification themselves involve further posits. This is the proposed
\emph{network of inductions}, not an assumption-free reduction of
all uncertainty to one numerical table.

The original melting example argues that unsuccessful attempts to
melt carbon need not license the final prediction that it will never
melt: other substances which initially failed to melt supply a
cross-inductive correction. This is the author's historical reasoning
example, not contemporary materials-science advice. In the historical
chapter's measurement example, row estimates \(0.8,0.7,0.79\)
are compared with \(0.8\pm0.05\); two of three lie in that band,
so the empirical second-level frequency is \(2/3\). That arithmetic
does not prove a \(2/3\) probability that an unknown first-level
limit lies in the band without a justified reference class and the
further assumptions of the appraisal.

The original states four grid convergence claims on pp.~342--343,
including finite-row attainability and an argument involving vertical
limits. Its claim of the network's superiority is distinct from those
conditional claims. Eberhardt and Glymour report Creary's later
criticisms of the grid/lattice choice, the competing correction rules,
the missing inference to faster convergence, and the restrictive
assumptions of the later theorems. The cited grid results and a full
mathematical reconstruction of their domains remain separate obligations;
we do not turn ``more rows'' into a proved uniform convergence rate.

The historical authors compare the network with hierarchical Bayesian
methods. This is an interpretive comparison, not the source's native
definition or an identity with a flat-prior update. A proper uniform
binary prior, for example, predicts \(2/3\) after one positive result;
the straight frequency guess is one. The related Carnap calculations
already proved in this book make the difference explicit. The historical
chapter itself questions how the straight rule and Bayesian-style
appraisal are combined and how their values become uniquely determined.

\section{The separate historical criticisms and their mathematical boundaries}
\label{sec:p15-reichenbach-historical-criticism}

The historical authors distinguish Popper's objections to probabilities
of scientific theories, Nagel's objections to the straight-rule
vindication, and Russell's objections to the logical foundation, while
not treating these as exclusively owned arguments. One proposed route
rates a theory by a sequence of testable consequences; another places
the theory in a reference class of theories. A success frequency among
consequences is not the truth of their unrestricted conjunction, and
a class of theories needs a criterion for membership and appraisal.
The source's discussions of measurement tolerance and schematized laws
are attempts to respond to this difference, not a proof that every
universal law is numerically identical to its confirmation ratio.

\begin{hbtheorem}{thm:p15-reichenbach-higher-order-boundary}{proved=annex:p15.reichenbach-higher-order-boundary}{Why a hierarchy is not a naive product}
If a higher-level claim \(H\) has probability \(q\) and an
event \(E\) has conditional degrees \(r\) on \(H\) and
\(s\) on its complement, then
\[
 P(E)=qr+(1-q)s.
\]
The product \(qr\) is the joint contribution from \(H\), not
the full event probability. With \(q=1/2,r=3/4,s=9/10\),
the full value is \(33/40\), not \(3/8\).
Nor does ``every factor is below one'' alone force an infinite product
to vanish: factors \(1-2^{-n-2}\), \(n\geq1\), have a positive
product. These results do not validate a particular hierarchy; they
identify assumptions needed by a multiplication objection.
\end{hbtheorem}

Reichenbach's response to Russell, reported through the historical
chapter, points to the missing probability of the event when the
higher-level proposition is false. This is distinct from multiplying
an infinite conjunction of confidence claims. The actual letter and
the later collected-writings version are not acquired originals here.
The criticisms by Lewis about uncertain epistemic foundations likewise
remain philosophical arguments, not automatic consequences of an
infinite product.

The historical chapter also reports Reichenbach's opposition to
Popper's verification/falsification asymmetry under uncertain empirical
measurements, his criticism of Carnap/Hempel-style choices of a
probability metric, and his preference for inverse probability over
maximum likelihood alone. Its raven discussion distinguishes a
probability implication from deductive contraposition and depends on
base frequencies. Its Goodman discussion questions whether the
``narrowest reference class'' response provides a principled choice
before the future cutoff. Neither discussion is credited as the
separate full Hempel, Popper, Carnap or Goodman presentation.

Finally, Eberhardt and Glymour separate uncertainty due to measurement
noise, underdetermination of a law by finite data, and Reichenbach's
interpretation of physical indeterminacy. They criticize representing
all three by one undifferentiated degree. Their comparison with Jeffrey
conditionalization is a later suggested framework, not a rule of the
original logic.

\begin{hbtheorem}{thm:p15-reichenbach-finite-fit}{proved=annex:p15.reichenbach-finite-fit}{The historical authors' noise-free fit failure}
Suppose finitely many distinct real inputs \(x_1,\ldots,x_m\)
have outputs on a line \(L(x)=ax+b\). For every real \(c\),
\[
 g_c(x)=L(x)+c\prod_{j=1}^m(x-x_j)
\]
fits all those observations. Distinct choices of \(c\) give different
values at every unsampled input. Thus even exact noise-free observations
do not determine that global functional form, absent further restrictions.
\end{hbtheorem}

This makes the historical chapter's final reasoning case explicit.
Adding other finite datasets changes the finite product but does not
remove this underdetermination by itself. A law-selection procedure
must declare its additional class restrictions or inductive commitments;
it cannot obtain them by calling the finite fit an objective frequency.

\begin{hbmethod}{def:p15-reichenbach-reasoning-method}{Using the source's method with its commitments visible}
Specify the event types, pairing, order, reference class and any phase
or normality assumptions. Keep actual finite counts separate from an
assertion about their continuation. Calculate relative and compound
degrees only where their values and denominators are defined; provide
the coupling rather than infer it from two marginals. State whether
the single-case commitment is an appraised or blind posit. When using
cross-inductions, identify the row/reference-class selection, the higher
level being estimated, the further probability assumptions and the
unappraised top level. Finally identify the predictive end: estimating
a limiting frequency, predicting an individual event, or choosing a
global law are different tasks, with different success conditions.
\end{hbmethod}

\begin{proofboundary}{def:p15-reichenbach-source-boundary}{What is reconstructed, and what remains open}
The annex proves the displayed counting realization, coupling table,
modal/threshold results and explicit convergence/reference-class
countercases. These are bounded formal results, not a theorem that
Reichenbach solved the problem of induction. The grid superiority,
complete normal-sequence construction, phase and continuous probability
sequence developments, and full syntax/semantics of iterated probability
implication are not supplied by these proofs.

The directly acquired 1937 French original and the English
\emph{Experience and Prediction} excerpt are not substituted for the
unacquired German 1935 or revised English 1949 books. The English excerpt
lacks its printing-identification leaf. The Caltech historical copy is
only thirty-two pages, printed 357--388: the allocated final bibliography
page 389 is missing, and its 2009 copyright footer/final2011 publisher
identity is not resolved. These are source dependencies, not grounds
for fabricating a complete presentation or whole-facet closure.
\end{proofboundary}
